% 第一章：绪论

\chapter{绪论}

\section{研究背景与意义}

近年来，以GPT系列\cite{brown2020gpt3,achiam2023gpt4}和Llama系列\cite{dubey2024llama3}为代表的大语言模型（Large Language Models，LLMs）在自然语言处理领域取得了突破性进展，在文本生成、机器翻译、问答系统、情感分析等多种任务上展现出卓越性能。这些模型凭借其强大的泛化能力，已在金融风控、医疗辅助、智能教育等安全敏感型行业实现了广泛部署，显著提升了相关领域的自动化水平和用户体验。

然而，大语言模型的广泛应用也带来了前所未有的安全挑战。\textbf{后门攻击}（Backdoor Attack）是一种尤为隐蔽且危险的训练期威胁。攻击者通过污染训练数据或直接修改模型参数，在模型中植入隐藏的触发器（Trigger），使模型在正常输入上表现如常，但一旦输入包含特定触发器，模型便会输出攻击者预设的恶意结果。后门攻击难以被检测的根本原因在于：触发器形式多样（可以是特定词语、句法结构、语义模式等），且对防御方而言通常是完全未知的；此外，大语言模型的庞大参数规模和复杂的参数化方式使得后门行为深埋于高维隐空间之中，难以被现有方法准确定位和消除。

随着大模型训练成本的急剧上升，模型的开发、微调乃至部署过程日益外包给第三方服务提供商，这使得模型在供应链各环节均面临被植入后门的风险。在此背景下，一种新型安全服务范式——\textbf{后门防御即服务}（Backdoor Defense-as-a-Service，BDaaS）应运而生。在这种范式下，模型所有者可在模型部署前将其提交至专业安全平台，由平台对模型进行后门净化，从而降低将受损模型投入生产的风险。这一设想与业界已有的商业化安全服务（如TrojAI Detect、CalypsoAI推理平台）的运作模式高度契合。

然而，现有后门防御方法在实际落地BDaaS场景时面临两大核心瓶颈。

\textbf{其一，对先验信息的强依赖。}许多方法需要目标任务的干净数据集、已知的触发器或目标输出信息，乃至任务领域信息。然而在通用大模型的服务化场景中，这些信息往往不可获取或涉及用户隐私，难以作为服务的预设条件。

\textbf{其二，缺乏可复用性。}现有方法通常针对每个受损模型独立进行净化，无法将之前积累的知识和防御组件复用到后续模型的净化中，严重制约了服务化平台的规模化能力和经济性。

鉴于上述挑战，本文提出ProtoPurify方法，通过在极少先验假设的条件下构建和利用后门原型向量，实现跨数据集、跨任务的通用后门净化，为大语言模型后门防御即服务提供了一种可行的技术路径。

\section{国内外研究现状}

\subsection{后门攻击研究}

后门攻击在深度学习领域已有较为深入的研究。根据触发器设计方式的不同，现有攻击方法大致分为三类：

\textbf{单触发器后门攻击}（Single-trigger Backdoor Attack）是最基础的攻击形式，使用单一固定触发器操控模型输出。Kurita等人\cite{kurita2020weight}针对BERT模型提出了基于稀有词的后门攻击；BadEdit\cite{li2024badedit}通过直接修改模型权重植入后门，无需污染训练数据。此类攻击依赖显式词汇线索，在一定程度上可被简单的异常词检测防御识别。

\textbf{多触发器后门攻击}（Multi-trigger Backdoor Attack）将触发器分散至多个输入部分，仅当所有触发器同时出现时才激活恶意行为，显著提高了攻击隐蔽性。复合后门攻击（CBA）\cite{huang2023cba}将触发器散布于用户查询与系统指令等不同提示片段；双触发器后门攻击\cite{hou2024dual}要求特定句法模板与虚拟语气模式同时出现方可激活，进一步增强了抗检测能力。

\textbf{无触发器后门攻击}（Triggerless Backdoor Attack）不依赖显式触发器字符串，而是利用语义或句法上的隐蔽条件激活恶意行为。虚拟提示注入（VPI）\cite{yan2024vpi}通过语义级触发器操控模型输出；分布式触发器后门攻击（DTBA）\cite{hao2024dtba}则将触发器分散于多轮对话之中，此类攻击对传统检测方法构成了更严峻的挑战。

\subsection{后门防御研究}

针对后门攻击的防御方法主要分为检测和净化两大方向。净化方法因能直接消除后门行为并恢复模型的可部署状态而更具实用价值，可进一步细分为以下四类：

\textbf{基于微调的方法}通过在干净数据上重新训练覆盖后门行为，代表方法包括NAD（神经注意力蒸馏）\cite{li2021nad}和W2SDefense\cite{yang2022w2sdefense}等。此类方法计算代价通常较高，且对高级后门攻击效果有限。

\textbf{基于剪枝的方法}通过剪除编码后门行为的参数子集实现净化，代表方法包括Fine-Pruning\cite{liu2018finepruning}、PURE\cite{zhu2024pure}和Obliviate\cite{li2024obliviate}等。这类方法依赖干净数据集识别需要剪除的参数，同样存在先验信息依赖问题。

\textbf{基于参数编辑的方法}将后门视为模型权重中的局部关联，通过直接修改小部分参数消除，代表方法包括ROME\cite{meng2022rome}和MEMIT\cite{meng2023memit}等。此类方法要求防御方知晓触发的不当行为，对未知或语义触发器效果有限。

\textbf{基于嵌入的方法}{\sloppy 
在表示空间而非直接在权重中运作，代表方法包括BEEAR\cite{zeng2024beear}和CROW\cite{min2024crow}等。BEEAR通过对抗性操控嵌入空间中的后门方向消除后门；CROW通过内部一致性正则化实现稳定的跨层表示。
\par}

上述方法虽在特定设置下取得了一定效果，但普遍存在对先验信息的强依赖或缺乏可复用机制等问题，难以直接适用于BDaaS场景。

\subsection{权重空间表示研究}

近期研究揭示了模型行为可以通过权重空间中的向量来捕获和操控这一重要规律。任务向量\cite{ilharco2023taskv}将模型在微调前后的参数差值形式化为编码特定任务能力的向量，这些向量具有近似线性的代数性质，可通过加减运算实现能力的组合或消除。模型合并（Model Soups）\cite{wortsman2022soups}等方法进一步证明，独立微调模型的权重平均可提升性能和鲁棒性。这些发现为本文提出的基于权重空间的后门净化方法提供了重要的理论支撑。

\section{主要研究内容与贡献}

本文围绕面向大语言模型后门防御即服务的核心需求，提出了基于原型表示的后门净化方法ProtoPurify，核心思想源于以下实证发现：不同后门攻击在模型参数空间中往往引发相关的参数变化，这意味着存在一个共享的恶意后门"原型"，可被识别和利用以实现通用净化。

本文的主要贡献如下：

\begin{enumerate}[leftmargin=2em]
    \item \textbf{提出ProtoPurify净化框架。}在极少先验假设（无需了解触发器、训练数据或下游任务信息）的条件下，通过构建原型向量实现后门净化，原生支持可复用性、可定制化、可解释性和运行时高效性，是面向BDaaS场景的实用防御框架。

    \item \textbf{构建端到端权重空间净化流程。}该流程包含四个核心阶段：通过配对模型模拟提取后门向量（阶段一）、通过聚合操作构建候选原型（阶段二）、通过逐层对齐分析检测边界层（阶段三），以及通过抑制原型对齐分量实现可控净化（阶段四）。

    \item \textbf{开展系统性实验验证。}在Llama3-8B和Mistral-7B两种主流大语言模型上，针对分类与生成任务进行广泛实验，结果表明ProtoPurify在面对单触发器、多触发器和无触发器后门时，净化效果与模型效用的综合表现始终优于6种代表性防御基线，同时对自适应后门攻击具有较强鲁棒性，且在无后门模型上不造成性能损失。
\end{enumerate}

\section{论文结构安排}

本文的结构安排如下：

\textbf{第一章}为绪论，介绍研究背景与意义、国内外研究现状、主要研究内容与贡献以及论文结构安排。

\textbf{第二章}为相关工作，系统梳理大语言模型的基础知识、后门攻击与防御的研究进展，以及权重空间表示的相关方法，为本文方法的提出奠定理论基础。

\textbf{第三章}为ProtoPurify方法，详细阐述威胁模型的定义、方法总体框架，以及后门模拟、原型构建、边界层检测和模型净化四个核心阶段的技术细节，并给出完整的算法伪代码。

\textbf{第四章}为实验与结论，介绍实验设置、与6种基线方法的全面对比结果、消融实验分析、面向BDaaS四大特性的分析、自适应攻击鲁棒性评估以及全文总结。
